\section*{Цель работы}
Приобретение практических навыков в управлении процессами в ОС и обеспечении обмена данными между процессами посредством каналов.

\section*{Задание}
Составить и отладить программу на языке Си, осуществляющую работу с процессами и взаимодействие между ними. В результате работы программа (основной процесс) должна создать для решения задачи один или несколько дочерних процессов. Взаимодействие между процессами осуществляется через системные сигналы/события и/или каналы (pipe). Необходимо обрабатывать системные ошибки.

\subsection*{Вариант 2}
Пользователь вводит команды вида: «число число число<endline>». Далее эти числа передаются от родительского процесса в дочерний. Дочерний процесс считает их сумму и выводит её в файл. Числа имеют тип float. Количество чисел может быть произвольным.